% ==================== PARTIE C — BOUCLES ====================
% Étape de conception : attente de verrouillage, série de sessions, recharge simulée, paliers.
%
% Progression :
%   C1  QCM — sémantique exacte de la borne d'arrêt
%   C2  corriger une boucle (compteur non incrémenté)
%   C3  trace d'une boucle de recharge à gain dégressif (piège : bornes >= atteintes pile à 60 et 80, arrêt pile à 90)
%   C4a compléter l'ossature d'une boucle (3 trous)
%   C4b prédire l'affichage dans un cas limite (0 tour)
%   C5  écrire le corps d'un main (bilan, maximum)

\element{boucle-qcm-borne}{
    \begin{question}{boucle-qcm-borne}
        Au branchement, la borne revérifie le verrouillage du câble toutes les 15 secondes
        tant que la durée d'attente écoulée ne dépasse pas 60 secondes.

        \begin{center}
            \begin{minipage}{.72\textwidth}
                \lstinputlisting[language=c]{sources/codes/c1-qcm-borne.c}
            \end{minipage}
        \end{center}

        Qu'affiche ce fragment ?
        \begin{multicols}{6}
            \begin{reponses}
                \bonne{\texttt{5 75}}
                \mauvaise{\texttt{4 60}}
                \mauvaise{\texttt{4 75}}
                \mauvaise{\texttt{5 60}}
                \mauvaise{\texttt{3 45}}
                \mauvaise{\texttt{4 45}}
            \end{reponses}
        \end{multicols}
    \end{question}
}

\element{boucle-correction}{
    \begin{question}{boucle-correction}
        Pour le bilan d'un point de charge, un développeur inexpérimenté a écrit le programme suivant, qui
        cumule l'énergie délivrée par chacune des sessions de la journée.

        \lstinputlisting[language=c]{sources/codes/c1b-moyenne-a-corriger.c}

        Ce programme ne se comporte pas comme prévu. Décrire le comportement observé à
        l'exécution, puis corriger le programme \textbf{directement sur le code}. Justifier
        ci-dessous.

        \evaluationProfGrille{2}{4}
    \end{question}
}

\element{boucle-trace}{
    \needspace{18\baselineskip}% évite de séparer l'intro du code
    Pour préserver la batterie, la borne réduit la puissance délivrée lorsque le niveau
    de charge augmente : le niveau gagne 20 points par pas de 10 minutes sous 60 \%,
    puis 10 points de 60 à 80 \%, puis 5 points au-delà. Le fragment ci-dessous simule
    cette recharge (\lstinline[language=C]|gain_pct| est le nombre de points de
    pourcentage gagnés pendant un pas de 10 minutes).
    
    \begin{question}{boucle-trace}
        Compléter le tableau en donnant, à la fin de \textbf{chaque tour de boucle}, les
        valeurs de \lstinline[language=C]|gain_pct|, \lstinline[language=C]|niveau_pct| et
        \lstinline[language=C]|minutes|. Numéroter vous-même les tours suivants et
        \textbf{s'arrêter au dernier tour effectivement exécuté} (les lignes en trop restent vides).
        \evaluationProf[2]

        \begin{minipage}[t]{.48\textwidth}
            \lstinputlisting[language=c]{sources/codes/c2-recharge-trace.c}
        \end{minipage}\hfill
        \begin{minipage}[t]{.48\textwidth}~\\
            {\renewcommand{\arraystretch}{2}
            \begin{tabular}{|c|p{1.95cm}|p{1.75cm}|p{1.35cm}|}
                \hline
                \textbf{Tour} & \lstinline[language=C]|gain_pct| & \lstinline[language=C]|niveau_pct| & \lstinline[language=C]|minutes| \\ \hline
                1 & & & \\ \hline
                2 & & & \\ \hline
                  & & & \\ \hline
                  & & & \\ \hline
                  & & & \\ \hline
                  & & & \\ \hline
                  & & & \\ \hline
            \end{tabular}}
        \end{minipage}
    \end{question}

    \begin{question}{boucle-trace-affichage}
        Qu'affiche ce programme après exécution complète de la boucle? 

        \evaluationProfGrille{2}{1}

    \end{question}
}

\element{boucle-completion}{
    \needspace{34\baselineskip}% évite de séparer l'intro du code
    \paragraph{Simulation de la montée en charge : }
    La borne monte le niveau de charge par paliers de 7 points de pourcentage, en partant
    de 15 \%. Le cahier des charges impose que :
    \begin{itemize}\setlength{\itemsep}{0pt}
        \item la boucle s'arrête dès que le niveau atteint \textbf{ou dépasse} le niveau
        cible saisi par l'utilisateur ;
        \item le niveau affiché ne dépasse \textbf{jamais} 100 \%.
    \end{itemize}

    \lstinputlisting[language=c,escapeinside={(*}{*)}]{sources/codes/c3-paliers-a-completer.c}
    
    \begin{question}{boucle-completion}
        Compléter les trois éléments manquants \textbf{directement dans les cadres pointillés}.

        \evaluationProf[2]
    \end{question}
    \begin{question}{boucle-completion-affichage}
        Indiquer ce que le programme affiche si l'utilisateur saisit un niveau cible de 10 \%.
    
        \evaluationProfGrille{2}{2}
    \end{question}
}


\element{boucle-conception}{
    \begin{question}{boucle-conception}
        L'exploitant veut un bilan de fin de journée pour un point de charge.

        Compléter le corps de la fonction \lstinline[language=C]|main| ci-dessous pour qu'elle :
        \begin{itemize}\setlength{\itemsep}{0pt}
            \item demande le nombre de sessions de la journée (entier) ;
            \item demande ensuite, une par une, l'énergie délivrée par chaque session (en kWh,
                  nombre à virgule) ;
            \item calcule l'énergie totale délivrée, le nombre de sessions ayant dépassé
                  20 kWh, et l'énergie de la session la plus importante de la journée ;
            \item affiche ces trois résultats à la fin.
        \end{itemize}

        \textbf{Contraintes :} les valeurs saisies ne doivent pas être conservées toutes en
        mémoire (on ne dispose d'aucune structure permettant de stocker une série de valeurs) ;
        chaque valeur est donc traitée au moment de sa saisie. Le choix de la boucle et la
        gestion correcte du maximum lors de la première itération font partie de l'évaluation.

        \lstinputlisting[language=c]{sources/codes/squelette-main-debut.c}
        \vspace*{-.5em}
        \evaluationProfGrille{4}{12}
        \lstinputlisting[language=c]{sources/codes/squelette-main-fin.c}
    \end{question}
}
